\section*{第 \ref{ch:b1:rock-forming-minerals} 章——造岩矿物}

\begin{solution}{exo:b1:rock-forming-minerals:1}
\ce{SiO4^{4-}}：\omterm{def:b1:rock-forming-minerals:silicate-classes}{岛状硅酸盐}，$s = 0$。\ce{Si2O7^{6-}}：\omterm{def:b1:rock-forming-minerals:silicate-classes}{双岛状硅酸盐}，
$s = 1$。\ce{Si6O18^{12-}}：\omterm{def:b1:rock-forming-minerals:silicate-classes}{环状硅酸盐}，$s = 2$。\ce{Si2O5^{2-}}：
\omterm{def:b1:rock-forming-minerals:silicate-classes}{层状硅酸盐}，$s = 3$。\ce{SiO2}：\omterm{def:b1:rock-forming-minerals:silicate-classes}{架状硅酸盐}，$s = 4$。\ce{SiO3^{2-}}
对应 $4 - s/2 = 3$，故 $s = 2$：可为环状（\omterm{def:b1:rock-forming-minerals:silicate-classes}{环状硅酸盐}）或单链
（\omterm{def:b1:rock-forming-minerals:silicate-classes}{链状硅酸盐}）；仅凭比值无法区分二者。
\end{solution}

\begin{solution}{exo:b1:rock-forming-minerals:2}
镁\omterm{def:b1:rock-forming-minerals:groups}{橄榄石}：$2\times24.305 + 28.085 + 4\times15.999 =
\qty{140.691}{g.mol^{-1}}$；铁\omterm{def:b1:rock-forming-minerals:groups}{橄榄石}：$2\times55.845 + 28.085 +
4\times15.999 = \qty{203.771}{g.mol^{-1}}$。摩尔体积：
$140.691/3.275 = \qty{42.96}{cm^3.mol^{-1}}$ 和 $203.771/4.392 =
\qty{46.40}{cm^3.mol^{-1}}$。
\end{solution}

\begin{solution}{exo:b1:rock-forming-minerals:3}
$(0.78 - 0.72)/0.72 = +\qty{8.3}{\percent}$，对 \ce{Fe^{2+}}；
$-\qty{4.2}{\percent}$，对 \ce{Ni^{2+}}；$+\qty{15.3}{\percent}$，对
\ce{Mn^{2+}}；$+\qty{38.9}{\percent}$，对 \ce{Ca^{2+}}。铁和镍，
均低于 \qty{15}{\percent}，可自由替代镁；略高于此值的锰，
替代受限；钙超过 \qty{30}{\percent}，几乎不能进入。四者
电荷均与镁相同。
\end{solution}

\begin{solution}{exo:b1:rock-forming-minerals:4}
$(3.51 - 2.16)/2.16 = \qty{62.5}{\percent}$，由石墨到金刚石；
$(2.93 - 2.71)/2.71 = \qty{8.1}{\percent}$，由方解石到文石。
密度更大的\omterm{def:b1:rock-forming-minerals:polymorph}{同质多象变体}，即金刚石和文石，是高压形式。
\end{solution}

\begin{solution}{exo:b1:rock-forming-minerals:5}
取 $x = 0.30$，代入 $\mathrm{Na}_{1-x}\mathrm{Ca}_x\mathrm{Al}_{1+x}\mathrm{Si}_{3-x}\mathrm{O}_8$：
$\mathrm{(Na_{0.7}Ca_{0.3})(Al_{1.3}Si_{2.7})O_8}$。电荷：
$0.7 + 2\times0.3 + 3\times1.3 + 4\times2.7 = 0.7 + 0.6 + 3.9 + 10.8 = +16$，
而 $8\times(-2) = -16$。
\end{solution}

\begin{solution}{exo:b1:rock-forming-minerals:6}
一个重复单元含两个四面体各共享三个角顶（其顶点
连接两条链），另两个各共享两个角顶。四个四面体共有
$4\times4 = 16$ 个角顶，其中 $3 + 3 + 2 + 2 = 10$ 个被共享，即
$10/2 = 5$ 个共享氧；其余 $16 - 10 = 6$ 个未共享。共计
$5 + 6 = 11$ 个氧对应 4 个硅：\ce{Si4O11}，电荷为
$4\times4 - 11\times2 = -6$。由命题取
$s = 5/2$ 得同样结果：每个硅对应 $4 - 5/4 = 11/4$ 个氧，每个硅电荷 $5/2 - 4 = -3/2$，
与上式一致。
\end{solution}

\begin{solution}{exo:b1:rock-forming-minerals:7}
\qty{100}{g} 中各组分的物质的量：\ce{SiO2} $40.20/60.083 = 0.66907$；\ce{MgO}
$47.10/40.304 = 1.16862$；\ce{FeO} $12.20/71.844 = 0.16981$；\ce{NiO}
$0.40/74.692 = 0.00536$。氧的物质的量：$2\times0.66907 + 1.16862 + 0.16981 +
0.00536 = 2.68194$；$f = 4/2.68194 = 1.49146$。按每四个氧计：Si
\num{0.998}，Mg \num{1.743}，Fe \num{0.253}，Ni \num{0.008}。
八面体间隙容纳 $1.743 + 0.253 + 0.008 = 2.004$，四面体间隙
\num{0.998}：分析结果良好，\ce{(Mg_{1.743}Fe_{0.253}Ni_{0.008})SiO4}。
$\mathrm{Fo} = 100\times1.743/(1.743 + 0.253) = 87.3$。因子 $f$
对 Mg 和 Fe 的作用相同，故在比值中相互抵消：
$\mathrm{Fo} = 100\,(47.10/40.304)/(47.10/40.304 + 12.20/71.844) =
100\times1.16862/1.33843 = 87.3$。
\end{solution}

\begin{solution}{exo:b1:rock-forming-minerals:8}
正电荷：$1 + 2\times3 + 3\times4 + 3 = +22$；负电荷：
$10\times(-2) + 2\times(-1) = -22$。若组成为 \ce{Si4O10} 的层，
应带 $-4$（见\cref{prop:b1:rock-forming-minerals:si-o}，$s = 3$）；将四个硅中的一个
换为铝，增加一个负电荷，故 \ce{(Si3Al)O10}
带 $-5$。两个八面体铝与两个羟基合计
$6 - 2 = +4$，平衡五个负电荷中的四个；层间的钾离子
平衡第五个。
\end{solution}

\begin{solution}{exo:b1:rock-forming-minerals:9}
(a) $M = (140.691 + 203.771)/2 = \qty{172.231}{g.mol^{-1}}$，$V = (42.96 +
46.40)/2 = \qty{44.68}{cm^3.mol^{-1}}$，$\rho = 172.231/44.68 =
\qty{3.855}{g.cm^{-3}}$；密度的平均值为 $(3.275 + 4.392)/2 =
\qty{3.834}{g.cm^{-3}}$，少 \qty{0.021}{g.cm^{-3}}。(b) 解方程
$203.771 + x(140.691 - 203.771) = 3.600\,\bigl(46.40 + x(42.96 - 46.40)\bigr)$：
$203.771 - 63.080x = 167.040 - 12.384x$，故 $x = 36.731/50.696 = 0.7245$：
\omterm{def:b1:rock-forming-minerals:groups}{橄榄石}为 $\text{Fo}_{72}$（精确到整数）。
\end{solution}

\begin{solution}{exo:b1:rock-forming-minerals:10}
其中 $M_i = \rho_i V_i$，
\[
  \rho(x) - \bigl(x\rho_1 + (1 - x)\rho_2\bigr)
  = \frac{x\rho_1V_1 + (1 - x)\rho_2V_2 - \bigl(x\rho_1 + (1 - x)\rho_2\bigr)\bigl(xV_1 + (1 - x)V_2\bigr)}{V(x)} .
\]
展开乘积，分子为 $x\rho_1V_1 - x^2\rho_1V_1 +
(1 - x)\rho_2V_2 - (1 - x)^2\rho_2V_2 - x(1 - x)(\rho_1V_2 + \rho_2V_1)$，而
$x - x^2 = x(1 - x)$，$(1 - x) - (1 - x)^2 = x(1 - x)$，故其为
$x(1 - x)(\rho_1V_1 + \rho_2V_2 - \rho_1V_2 - \rho_2V_1) =
x(1 - x)(\rho_1 - \rho_2)(V_1 - V_2)$。对\omterm{def:b1:rock-forming-minerals:groups}{橄榄石}（1 为镁\omterm{def:b1:rock-forming-minerals:groups}{橄榄石}），两个
差值均为负，$3.275 - 4.392 = -1.117$ 和 $42.96 - 46.40 =
-3.44$，故二者之积为正，密度位于直线
上方。在 $\text{Fo}_{50}$ 处：$0.25\times1.117\times3.44/44.68 =
\qty{0.0215}{g.cm^{-3}}$，与\cref{exo:b1:rock-forming-minerals:9}的结果
在摩尔体积取舍的误差范围内相符。
\end{solution}

\begin{solution}{exo:b1:rock-forming-minerals:11}
总量 \qty{99.60}{\percent}。各氧化物的物质的量：\ce{SiO2} $0.90042$，
\ce{Al2O3} $0.28344$，\ce{CaO} $0.19973$，\ce{Na2O} $0.08229$，
\ce{K2O} $0.00318$。氧的物质的量：$1.80084 + 0.85033 + 0.19973 + 0.08229 +
0.00318 = 2.93636$；$f = 8/2.93636 = 2.72446$。阳离子（对
\ce{Al2O3}、\ce{Na2O}、\ce{K2O} 取两倍物质的量）：Si \num{2.453}，Al \num{1.544}，Ca
\num{0.544}，Na \num{0.448}，K \num{0.017}。四面体间隙：
$2.453 + 1.544 = 3.997$；大阳离子位：$0.544 + 0.448 + 0.017 = 1.009$。
于是 $\mathrm{Al} - 1 = 0.544$，在舍入范围内与 Ca 相等：每个钙都与一个铝一同进入。
$\mathrm{An} = 100\times0.544/1.009 = 53.9$：属
拉\omterm{def:b1:rock-forming-minerals:groups}{长石}，$\text{An}_{54}$。
\end{solution}

\begin{solution}{exo:b1:rock-forming-minerals:12}
(a) $M(\ce{Fe3O4}) = 3\times55.845 + 4\times15.999 = \qty{231.531}{g.mol^{-1}}$.
其三个铁以 \ce{FeO} 计，质量为 $3\times71.844 = \qty{215.532}{g}$，
即 $215.532/231.531 = \qty{93.09}{\percent}$；与
\ce{Fe^{3+}} 结合的氧（超出 \ce{FeO} 所能计入的部分）未计入总量。
(b) \ce{FeO} 与氧的物质的量均为 $93.09/71.844 = 1.29572$；$f =
4/1.29572 = 3.08708$；四个铁对应四个氧，即 \ce{Fe4O4}。磁铁矿
中三个阳离子对应四个氧。氧的数目被低估：每个
\ce{Fe^{3+}} 结合一个半氧，而非 \ce{FeO} 所假设的一个。(c) 以三个阳离子为基准，Fe $= 3$
设有 $a$ 个 \ce{Fe^{2+}} 与 $b$ 个 \ce{Fe^{3+}}：$a + b = 3$，$2a + 3b = 8$，
于是 $b = 2$，$a = 1$：$\mathrm{Fe^{2+}Fe^{3+}_2O_4}$。(d) 每个化学式单元含一个 \ce{FeO} 与
一个 \ce{Fe2O3}：\ce{FeO} 占 $71.844/231.531 = \qty{31.03}{\percent}$
，\ce{Fe2O3} 占 $159.687/231.531 = \qty{68.97}{\percent}$，合计
\qty{100.00}{\percent}.
\end{solution}

\begin{solution}{pb:b1:rock-forming-minerals:1}
\textbf{1.} \ce{SiO2} $28.085 + 2\times15.999 = 60.083$; \ce{FeO} 71.844;
\ce{MnO} 70.937; \ce{MgO} 40.304; \ce{NiO} 74.692; \ce{CaO} 56.077
(\si{g.mol^{-1}}).

\textbf{2.} $39.50 + 16.53 + 0.23 + 43.32 + 0.20 + 0.22 =
\qty{100.00}{\percent}$。总量为 \qty{97}{\percent}，提示存在
一个未分析的组分（水，或与 \ce{Fe^{3+}} 结合而以
\ce{FeO} 形式记入的氧），或分析质量较差。

\textbf{3.} \ce{SiO2} $39.50/60.083 = 0.65742$; \ce{FeO} 0.23008; \ce{MnO}
0.00324; \ce{MgO} 1.07483; \ce{NiO} 0.00268; \ce{CaO} 0.00392.

\textbf{4.} $2\times0.65742 = 1.31485$ 为 \ce{SiO2} 的值，其余与
第 3 题相同；总和为 $1.31485 + 0.23008 + 0.00324 + 1.07483 +
0.00268 + 0.00392 = 2.62960$。

\textbf{5.} 每种氧化物含一个阳离子：数值同第 3 题。

\textbf{6.} $f = 4/2.62960 = 1.52114$.

\textbf{7.} Si $0.65742\times1.52114 = 1.000$; Fe 0.350; Mn 0.005; Mg 1.635;
Ni 0.004; Ca 0.006.

\textbf{8.} 八面体间隙之和 $1.635 + 0.350 + 0.005 + 0.004 + 0.006 = 2.000$，
四面体间隙为 1.000：恰为理想值 2 和
1.

\textbf{9.} \ce{(Mg_{1.635}Fe_{0.350}Mn_{0.005}Ni_{0.004}Ca_{0.006})SiO4}.

\textbf{10.} $2\times2.000 + 4\times1.000 = +8.000$，与 $4\times(-2) =
-8$。每个阳离子按其氧化物的电荷计入，因此
正电荷总数按构造恰为氧数的两倍，此检验不可能失败。
第 8 题的位点总数则可能失败。

\textbf{11.} 镍与镁的差异为
$(0.69 - 0.72)/0.72 = -\qty{4}{\percent}$，远低于 \qty{15}{\percent}：
故可自由替代。钙则偏大 $(1.00 - 0.72)/0.72 = \qty{39}{\percent}$
，超过 \qty{30}{\percent}：几乎不进入。

\textbf{12.} 将铁按 \ce{FeO} 计入而其中一部分为 \ce{Fe^{3+}} 时，
每个阳离子对应的氧过少，因此四个氧上的阳离子多于三个
（见\cref{exo:b1:rock-forming-minerals:12}）。总和为 $1.000 + 2.000 = 3.000$：不含显著的
\ce{Fe^{3+}}.

\textbf{13.} $\mathrm{Fo} = 100\times1.635/(1.635 + 0.350) = 163.5/1.985 =
82.4$.

\textbf{14.} $\mathrm{Mg\#} = 82.4$：所有铁均为 \ce{Fe^{2+}}
（第 12 题），因此 $\mathrm{Fe^{2+}} = \mathrm{Fe}$，两种定义
重合。

\textbf{15.} $\text{Fo}_{82.4}\text{Fa}_{17.6}$，$x = 0.824$；铁占
\qty{17.6}{\percent} 的八面体间隙（二元模型中）
（计入次要元素时为 $0.350/2.000 = \qty{17.5}{\percent}$）。

\textbf{16.} 因子同乘于 Mg 与 Fe，因而约去：
$\mathrm{Fo} = 100\times1.07483/(1.07483 + 0.23008) = 82.4$.

\textbf{17.} $40.5/40.304 = 1.00486$ 与 $20.0/71.844 = 0.27838$；
$\mathrm{Fo} = 100\times1.00486/1.28324 = 78.3$。较晚的\omterm{def:b1:rock-forming-minerals:groups}{橄榄石}
富铁。

\textbf{18.} $M = 140.691$ 与 \qty{203.771}{g.mol^{-1}}；
$V = 140.691/3.275 = 42.96$ 与 $203.771/4.392 =
\qty{46.40}{cm^3.mol^{-1}}$.

\textbf{19.} $M = 0.824\times140.691 + 0.176\times203.771 = 115.929 +
35.864 = \qty{151.793}{g.mol^{-1}}$; $V = 0.824\times42.96 +
0.176\times46.40 = 35.399 + 8.166 = \qty{43.565}{cm^3.mol^{-1}}$.

\textbf{20.} $\rho = 151.793/43.565 = \qty{3.484}{g.cm^{-3}}$.

\textbf{21.} $0.824\times3.275 + 0.176\times4.392 = 2.699 + 0.773 =
\qty{3.472}{g.cm^{-3}}$，\qty{0.012}{g.cm^{-3}}（\qty{0.3}{\percent}）低于
理想混合值：接近，正如 \cref{exo:b1:rock-forming-minerals:10}
所预测。

\textbf{22.} $0.12\times\qty{3484}{kg.m^{-3}}\times\qty{1}{m^3} =
\qty{418}{kg}$.

\textbf{23.} $418/2950 = \qty{14.2}{\percent}$：致密\omterm{def:b1:rock-forming-minerals:groups}{橄榄石}在质量中所占
比例高于其在体积中所占比例。

\textbf{24.} 该\omterm{def:b1:rock-forming-minerals:groups}{橄榄石}为 $\text{Fo}_{82}$：其镁\omterm{def:b1:rock-forming-minerals:groups}{橄榄石}含量为
\textbf{82\%}。
\end{solution}
